\emph{Non-identifiability of the moment-type diagonal coordinate from the Gram spectrum.} Within the named constant-variance gaussian edge-block subclass of the fourth preceding paragraph, at \(D=2\), write \(G=AA^{\dagger}\), \(\tau=\operatorname{Tr}G\), the Gram-spectrum coordinate \(p=(\lambda_{1},\lambda_{2})/\tau\) of that paragraph, and the diagonal statistic \(\hat\rho_{2}=\sum_{i}(G_{ii}/\tau)^{2}\), the second power-sum of the normalized diagonal profile \(d_{i}=G_{ii}/\tau\), which differs by the constant shift one from the moment-map coordinate \(\rho_{2}=1+s_{2}/W^{2}\) of the third preceding paragraph evaluated at \(d\), and which in the rank-one edge-datum class of the fourth preceding paragraph is exactly the slot-power profile statistic of that paragraph. The Schur--Horn fibre of the spectrum map is explicit: for a fixed spectrum \(\lambda_{1}\ge\lambda_{2}>0\) the entry \(G_{11}\) ranges exactly over the interval from \(\lambda_{2}\) to \(\lambda_{1}\), with \(G_{22}=\tau-G_{11}\) and \(\lvert G_{12}\rvert^{2}=G_{11}G_{22}-\lambda_{1}\lambda_{2}\), and every value in the interval is realized by an edge-block configuration of the subclass; consequently \(\hat\rho_{2}\) takes, on the fibre over any single nondegenerate spectrum value, exactly the interval from \(1/2\) to \((\lambda_{1}^{2}+\lambda_{2}^{2})/\tau^{2}\), of width \((\lambda_{1}-\lambda_{2})^{2}/(2\tau^{2})\), while at the fully degenerate spectrum \(\lambda_{1}=\lambda_{2}\) the fibre collapses to the scalar multiple of the identity and \(\hat\rho_{2}\) is forced to \(1/2\). Two configurations of trace \(4\) exhibit the failure: the diagonal Gram \(\operatorname{diag}(3,1)\) and the Gram with diagonal \((2,2)\) and \(\lvert G_{12}\rvert=1\) share the spectrum \((3,1)\) while \(\hat\rho_{2}\) takes the distinct values \(5/8\) and \(1/2\). In the rank-one edge-datum class of the fourth preceding paragraph the spectrum coordinate is the single atom at the vertex, identical for every configuration, while \(\hat\rho_{2}\) ranges over the interval from \(1/2\) to \(1\), the endpoints realized by the configurations of squared norm sequence \((1,0)\) and \((1/2,1/2)\). There therefore exists no map from the profile simplex to the real line -- measurable, relabelling-equivariant, or otherwise -- that recovers \(\hat\rho_{2}\) from the Gram-spectrum coordinate almost everywhere on the support of the law of any constant-variance gaussian edge-block datum at \(D=2\): the spectrum map forgets the diagonal completely off the fully degenerate spectrum, and this is a proved nonexistence on the support for these two named coordinates within the named subclass, not an unproved absence of a supplied formula; the identification holds only after an extra selection, for instance on diagonal Gram matrices, where spectrum and diagonal agree as multisets. The normalized-diagonal coordinate of the sixth preceding paragraph remains the carrier of the profile information within these named classes. These statements are exact linear-algebraic identifiability statements on the support of the named subclass at \(D=2\), they bound nonexistence between the two named coordinates only, they evaluate no matrix integral of the core action, they supply no state and no selection mechanism, they do not fix the state that sets \(W\), \(g\), \(s\) or the frame, they leave the missing steps (i)-(iv) of the Missing step paragraph unchanged, and they close nothing in Problem~G.
